\section{Conclusion}

Quantum game theory extends classical strategic interaction into the quantum domain by enabling players to manipulate entangled quantum states through local unitary operations. Within the EWL framework, entanglement alters equilibrium structure and reveals cooperative quantum equilibria that do not emerge in the classical prisoner's dilemma.

This work suggests that quantum equilibrium search may be interpreted as optimization over continuous parameterized unitary strategy spaces. Gradient-based interpretations of strategy adaptation suggest parallels with reinforcement learning and policy optimization.

Simultaneously, equilibrium stability is governed by the admissible strategy space. Restricted strategy classes enable cooperative equilibria, whereas expanded unitary spaces may introduce destabilizing counter-strategies.

As quantum computation and quantum machine learning continue to develop, quantum game-theoretic models may provide useful frameworks for studying strategic adaptation in computationally complex environments.